\pdfvariable suppressoptionalinfo 611
\providecommand{\CRevisionRound}{2}
\documentclass[11pt]{article}
\usepackage[a4paper,margin=22mm]{geometry}
\usepackage{amsmath,amssymb,amsthm,booktabs,microtype}
\usepackage[T1]{fontenc}
\usepackage{lmodern}
\usepackage[hidelinks,hypertexnames=false]{hyperref}
\emergencystretch=2em
\newtheorem{theorem}{Theorem}
\newtheorem{proposition}{Proposition}
\newtheorem{lemma}{Lemma}
\title{Energy Topology, Action, and the Homoclinic Boundary of the Duffing Flow}
\author{HCS Research Program}
\date{29 August 2026\quad (revision \CRevisionRound)}
\begin{document}
\maketitle

\begin{abstract}
We close an all-parameter phase-space atlas for the conservative Duffing
oscillator $\dot x=v$, $\dot v=-\delta x-\beta x^3$ with $\beta>0$.
The energy levels are classified for both the single-well and double-well
regimes, including the two homoclinic loops at the barrier.  Turning points
are roots of an explicit quadratic in $x^2$; an endpoint-cancelled integral
therefore gives the period and action on every regular oval, with the action
derivative identity and center, quartic, and saddle limits.  The ledger keeps
the $\beta=0$ faces separate and treats the continuum of energy ovals as a
source Hamiltonian family, not as isolated primitive cycles.  Exact symbolic,
independent numerical, replay, and hostile-mutation checks accompany the
proof.
\end{abstract}

\section{Model and invariant}
Put
\begin{equation}
 \dot x=v,\qquad \dot v=-\delta x-\beta x^3,
 \qquad H(x,v)=\frac{v^2}{2}+V(x),\quad
 V(x)=\frac{\delta x^2}{2}+\frac{\beta x^4}{4},
 \label{eq:model}
\end{equation}
where $\beta>0$ and $\delta\in\mathbb R$.  The physical time is the common
clock.  Hamiltonian conservation follows immediately from
$H_x\dot x+H_v\dot v=0$.  For a regular connected oval with turning points
$\ell<r$, define
\begin{equation}
 T(E)=\sqrt{2}\int_\ell^r\frac{dx}{\sqrt{E-V(x)}},\qquad
 I(E)=\frac1\pi\int_\ell^r\sqrt{2(E-V(x))}\,dx .
 \label{eq:quadrature}
\end{equation}
The endpoint square roots are integrable; the executable ledger evaluates
\eqref{eq:quadrature} after $x=(\ell+r)/2+(r-\ell)\sin\theta/2$.

\begin{theorem}[complete Duffing energy atlas]\label{thm:atlas}
Let $\beta>0$.  If $\delta\geq0$, every $E>0$ is one smooth compact oval
around the origin.  If $\delta<0$, write $a=\sqrt{-\delta/\beta}$ and
$V_{\min}=-\delta^2/(4\beta)$.  Then $V_{\min}<E<0$ gives two ovals,
$E=0$ gives two homoclinic loops to $(0,0)$, and $E>0$ gives one outer oval.
The squared turning points are
\begin{equation}
 y_\pm=\frac{-\delta\pm\sqrt{\delta^2+4\beta E}}{\beta},\qquad y=x^2 .
 \label{eq:roots}
\end{equation}
On each regular component, \eqref{eq:quadrature} is finite, $I'(E)=T(E)/(2\pi)$,
and the period is positive.  For $\delta=-\alpha^2<0$, the separatrices are
\begin{equation}
 x_h(t)=\pm\frac{\sqrt{2}\,\alpha}{\sqrt\beta}\,\operatorname{sech}(\alpha t),
 \qquad v_h=\dot x_h .
 \label{eq:homoclinic}
\end{equation}
Moreover $T(E)\to2\pi/\sqrt\delta$ as $E\downarrow0$ when $\delta>0$,
$T(E)\,E^{1/4}\beta^{1/4}$ is constant when $\delta=0$, and $T(E)$
diverges logarithmically as a regular energy approaches the saddle level.
\end{theorem}

\begin{proof}
The quartic potential is coercive, so regular compact components are the
connected components described by the signs of the two roots in
\eqref{eq:roots}.  For $\delta<0$ the identity
\[
 V(x)+\frac{\delta^2}{4\beta}=\frac\beta4\left(x^2+\frac\delta\beta\right)^2
\]
locates the two minima and the barrier.  Separation of variables gives
\eqref{eq:quadrature}; the sine substitution removes endpoint singularities.
Differentiating the action on a compact subinterval of regular energies and
using the vanishing endpoint integrand gives $I'=T/(2\pi)$.  At $\delta=0$,
the scaling $x=E^{1/4}\beta^{-1/4}u$ proves the quartic law.  The quadratic
Taylor expansion at a nondegenerate center gives the center limit.  Near the
hyperbolic saddle, $E-V(x)$ has a nondegenerate quadratic zero, and comparison
with $\int dx/\sqrt{\alpha^2x^2+|E|}$ gives the logarithmic divergence.
Finally direct substitution of \eqref{eq:homoclinic} into
$\ddot x-\alpha^2x+\beta x^3=0$ proves the separatrix formula.
\end{proof}

\section{Linear faces and source boundaries}
The origin has linear rates $\pm i\sqrt\delta$ for $\delta>0$ and
$\pm\sqrt{-\delta}$ for $\delta<0$; the well equilibria $x=\pm a$ have
rates $\pm i\sqrt{-2\delta}$.  When $\beta=0$, the positive-$\delta$ face is
the harmonic oscillator, the negative-$\delta$ face is inverted and has no
compact positive-energy ovals, and $(\beta,\delta)=(0,0)$ is free motion.
These lower-dimensional models are not silently folded into the quartic
theorem.  The two inner components and the outer component are separated in
the ledger, so a period near the barrier cannot be mistaken for a center
period.

\ifnum\CRevisionRound>0
\section{Action coordinates and evidence}
On each regular component the action is strictly increasing because
$I'(E)>0$; it is therefore a valid local action coordinate, while the
homoclinic level is its logarithmic endpoint.  The certificate contains five
parameter cases and twenty energy rows, including both double-well components
and the pure-quartic scaling control.  A producer-independent checker makes
790 assertions, and SymPy verifies ten identities (conservation, factorization,
turning polynomial, the profile residual, linear rates, and the action
integrand).  Clean replay is byte-for-byte and the hostile suite rejects
twenty repaired-hash/schema mutations plus one stale-hash mutation
(21/21 hostile checks).
\fi

\ifnum\CRevisionRound>1
\section{Route-A boundary}
The Duffing flow is a natural one-degree-of-freedom Hamiltonian system, but
its recurrent ovals are indexed by a continuum of real energies.  Thus the
source theorem supplies no intrinsic isolated primitive-orbit owner, logarithmic
prime clock, or target divisor.  The strict tuple is
\[
 (\texttt{A0\_FAIL,A1\_WEAK,A2\_FAIL,A3\_FAIL,A4\_FORMAL\_HINT}),
\]
with \texttt{ROUTE\_A\_REJECTED} and
\texttt{route\_b\_invocation\_allowed=false}.  The action variable and a
possible semiclassical lift are candidate-local formal hints only.  No
arithmetic target is inferred from the exact elliptic/quadrature structure.
\fi

\section*{Source note and declarations}
The phase-plane terminology follows Guckenheimer and Holmes~\cite{GH1983};
period-function conventions follow Chicone~\cite{Chicone2006}.  The name
Duffing is historical~\cite{Duffing1918}; no priority claim is made.  No
prime or zero table, local arithmetic datum, Euler factor, root number,
automorphy object, target functional equation, Hilbert--Pólya operator or
Route-B input is used.  Generative tools assisted drafting and code generation;
the displayed claims are tied to the artifact chain.  This is not external
peer review.  Scope literal: \texttt{NO\_BAD\_EULER\_OR\_ROOT\_NUMBER}.

\begin{thebibliography}{9}
\bibitem{GH1983} J. Guckenheimer and P. Holmes,
\emph{Nonlinear Oscillations, Dynamical Systems, and Bifurcations of Vector
Fields}, Springer, 1983. DOI: 10.1007/978-1-4612-1140-2.
\bibitem{Chicone2006} C. Chicone, \emph{Ordinary Differential Equations with
Applications}, 2nd ed., Springer, 2006. DOI: 10.1007/0-387-35794-7.
\bibitem{Duffing1918} G. Duffing, \emph{Erzwungene Schwingungen bei
veränderlicher Eigenfrequenz}, Vieweg, Braunschweig, 1918.
\end{thebibliography}
\end{document}
